\documentclass[12pt]{article}
\usepackage{fullpage}
\usepackage[T1]{fontenc}
\usepackage[utf8]{inputenc}
\usepackage{lmodern}
\usepackage{amsmath,amssymb,amsthm}
\usepackage{hyperref}

\newtheorem{theorem}{Theorem}
\newtheorem{lemma}{Lemma}
\newtheorem{proposition}{Proposition}
\newtheorem{corollary}{Corollary}
\theoremstyle{remark}
\newtheorem*{remark}{Remark}

\DeclareMathOperator{\score}{score}
\newcommand{\bR}{\mathbb{R}}
\newcommand{\boxp}{\boxplus_n}
\newcommand{\inner}[2]{\langle #1,#2\rangle}

\title{Subharmonicity of $1/\Phi_n$ under finite free convolution:\\
the Gaussian case, an exact de Bruijn identity, and an honest status}
\author{}
\date{}

\begin{document}
\maketitle

\begin{abstract}
We study the conjectural finite free analogue of Stam's inequality,
\[
\frac{1}{\Phi_n(p\boxp q)}\ \ge\ \frac{1}{\Phi_n(p)}+\frac{1}{\Phi_n(q)},
\qquad p,q\ \text{monic, real-rooted, degree }n. \tag{$\star$}
\]
We establish the following \emph{rigorously and unconditionally}: a closed form
for $\Phi_n$; covariance and the identity element; the finite free Gaussian
(Hermite) heat semigroup and the fact that convolution with it is the heat flow
$\partial_t p=-\tfrac12 p''$; an exact de~Bruijn identity
$\frac{d}{dt}\Phi_n(p_t)=-I(p_t)\le0$ with explicit nonnegative dissipation
$I$; the case $n=2$ with equality; and -- the main new contribution -- a sharp
\emph{finite free Cram\'er--Rao / Stam lower bound}
\[
I(p)\ \ge\ \frac{4}{n(n-1)}\,\Phi_n(p)^2,\qquad\text{equivalently}\qquad
g(p):=\frac{I(p)}{\Phi_n(p)^2}\ \ge\ \frac{4}{n(n-1)},
\]
with equality exactly at the Hermite polynomial.  From this we deduce a
complete proof of $(\star)$ \emph{whenever one of the two factors is a (scaled,
translated) finite free Gaussian} $G_t$, for every $n$, with explicit (strict
unless $p$ is Hermite) slack.  Finally we give an honest account of the
general case: we show that the natural ``infinitesimal Stam'' inequality along
a coupled heat flow is \emph{infeasible at genuinely reachable states}, so that
heat-flow coupling alone cannot prove $(\star)$ in general, and we isolate the
exact remaining open ingredient (a finite free score-projection lemma).
\end{abstract}

\section{Setup and notation}

Throughout, $p(x)=\prod_{i=1}^n(x-\lambda_i)$ is monic of degree $n$ with real
roots $\lambda_1,\dots,\lambda_n$.  When the $\lambda_i$ are distinct we write
\[
s_i \;=\; s_i(p)\;:=\;\sum_{j\neq i}\frac{1}{\lambda_i-\lambda_j},
\qquad i=1,\dots,n,
\]
for the components of the finite free score vector $\score(\lambda)$, and
\[
\Phi_n(p)\;=\;\|\score(\lambda)\|^2\;=\;\sum_{i=1}^n s_i^2 .
\]
If $p$ has a repeated root we set $\Phi_n(p)=+\infty$, so that $1/\Phi_n(p)=0$;
the degenerate cases reduce to the simple-root case by continuity
(Section~\ref{sec:degenerate}), and we assume below that all polynomials have
simple roots.

The finite free convolution is the bilinear operation on coefficients given in
the problem statement.  It is commutative and associative (both immediate from
the symmetric bilinear coefficient formula), with identity $x^n$.

\section{A closed form for $\Phi_n$ and the score Laplacian}\label{sec:identity}

\begin{lemma}\label{lem:phiform}
For $p$ with distinct roots,
\[
\Phi_n(p)\;=\;\sum_{i=1}^n\Big(\sum_{j\neq i}\frac{1}{\lambda_i-\lambda_j}\Big)^2
\;=\;\sum_{i\neq j}\frac{1}{(\lambda_i-\lambda_j)^2}.
\]
\end{lemma}

\begin{proof}
Expanding the square,
\[
\Phi_n(p)=\sum_{i=1}^n\sum_{j\neq i}\sum_{k\neq i}
\frac{1}{(\lambda_i-\lambda_j)(\lambda_i-\lambda_k)}
=\underbrace{\sum_{i\neq j}\frac{1}{(\lambda_i-\lambda_j)^2}}_{\text{terms }j=k}
\;+\;T,
\]
where $T=\sum_{i,j,k\ \text{distinct}}
\frac{1}{(\lambda_i-\lambda_j)(\lambda_i-\lambda_k)}$.
Fix distinct $a=\lambda_i,b=\lambda_j,c=\lambda_k$.  The partial-fraction
identity
\[
\frac{1}{(a-b)(a-c)}+\frac{1}{(b-a)(b-c)}+\frac{1}{(c-a)(c-b)}=0
\]
holds (clearing denominators gives $(b-c)+(c-a)+(a-b)=0$).  In $T$ each
unordered triple $\{i,j,k\}$ contributes exactly these three terms (each
element playing the centre role once); by the identity they cancel, so $T=0$.
\end{proof}

We now introduce the operator that organizes all the differential identities
below.  For a vector $v\in\bR^n$ define the \emph{score Laplacian}
\[
(Lv)_i\;:=\;\sum_{j\neq i}\frac{v_i-v_j}{(\lambda_i-\lambda_j)^2},
\qquad i=1,\dots,n .
\]
$L$ is the weighted graph Laplacian on the complete graph with edge weights
$w_{ij}=(\lambda_i-\lambda_j)^{-2}>0$; hence $L$ is symmetric and positive
semidefinite, with
\[
\inner{v}{Lv}=\tfrac12\sum_{i\neq j}\frac{(v_i-v_j)^2}{(\lambda_i-\lambda_j)^2}\ \ge\ 0,
\qquad \ker L=\operatorname{span}\{(1,\dots,1)\}. \tag{2.1}
\]

\begin{lemma}[Score-Laplacian identities]\label{lem:Lid}
Write $\lambda=(\lambda_1,\dots,\lambda_n)$ and $s=\score(\lambda)$.  Then
\begin{enumerate}
\item[(i)] $L\lambda=s$;
\item[(ii)] $\Phi_n(p)=\inner{s}{s}=\inner{L\lambda}{s}=\inner{\lambda}{Ls}$;
\item[(iii)] $\inner{\lambda}{s}=\inner{\lambda}{L\lambda}=\dfrac{n(n-1)}{2}$;
\item[(iv)] the de~Bruijn dissipation (Theorem~\ref{thm:debruijn}) equals
$I(p)=2\inner{s}{Ls}$.
\end{enumerate}
\end{lemma}

\begin{proof}
(i) $(L\lambda)_i=\sum_{j\neq i}\frac{\lambda_i-\lambda_j}{(\lambda_i-\lambda_j)^2}
=\sum_{j\neq i}\frac{1}{\lambda_i-\lambda_j}=s_i$.

(ii) $\Phi_n=\sum_i s_i^2=\inner{s}{s}$, and $\inner{s}{s}=\inner{L\lambda}{s}$
by (i).

(iii) By (i), $\inner{\lambda}{L\lambda}=\inner{\lambda}{s}
=\sum_i\lambda_is_i=\sum_{i\neq j}\frac{\lambda_i}{\lambda_i-\lambda_j}$.
Pairing the terms $(i,j)$ and $(j,i)$,
\[
\frac{\lambda_i}{\lambda_i-\lambda_j}+\frac{\lambda_j}{\lambda_j-\lambda_i}
=\frac{\lambda_i-\lambda_j}{\lambda_i-\lambda_j}=1,
\]
and there are $\binom n2$ unordered pairs, so the sum is $\binom n2=n(n-1)/2$.

(iv) is the definition of $I$ together with (2.1) applied to $v=s$:
$2\inner{s}{Ls}=\sum_{i\neq j}\frac{(s_i-s_j)^2}{(\lambda_i-\lambda_j)^2}=I(p)$.
\end{proof}

\section{Covariance and the identity element}\label{sec:cov}

\begin{lemma}\label{lem:cov}
Let $p,q$ be monic of degree $n$, $c\in\bR$, $s>0$.
\begin{enumerate}
\item[(i)] $p\boxp x^n=p$.
\item[(ii)] $p\boxp (x-c)^n=p(x-c)$; hence $\score$ and $\Phi_n$ are
translation invariant, and $(\star)$ is translation invariant in each argument.
\item[(iii)] If $p_s(x):=s^{n}p(x/s)$ is the polynomial with roots $s\lambda_i$,
then $\score(s\lambda)=s^{-1}\score(\lambda)$, $\Phi_n(p_s)=s^{-2}\Phi_n(p)$,
and $p_s\boxp q_s=(p\boxp q)_s$.
\end{enumerate}
\end{lemma}

\begin{proof}
Write $a_k,b_k,c_k$ for the coefficients of $p,q,p\boxp q$.

(i) For $q=x^n$, $b_0=1$, $b_j=0$ $(j\ge1)$; the only surviving term in
$c_k=\sum_{i+j=k}\frac{(n-i)!(n-j)!}{n!(n-k)!}a_ib_j$ is $j=0,i=k$, giving
$c_k=a_k$.

(ii) $(x-c)^n$ has $b_j=\binom nj(-c)^j$. With $i=k-j$ and $n-i=n-k+j$,
\[
\frac{(n-i)!(n-j)!}{n!(n-k)!}\binom nj
=\frac{(n-k+j)!\,(n-j)!}{n!\,(n-k)!}\cdot\frac{n!}{j!(n-j)!}
=\frac{(n-k+j)!}{(n-k)!\,j!}
=\binom{\,n-k+j\,}{j}.
\]
Hence
\[
c_k=\sum_{j=0}^{k}\binom{\,n-k+j\,}{j}\,a_{k-j}\,(-c)^j .
\]
We claim this equals $\sum_{j=0}^k\binom kj a_{k-j}(-c)^j$, the coefficient of
$x^{n-k}$ in $p(x-c)$.  Both polynomials $p\boxp(x-c)^n$ and $p(x-c)$ are
$\bR$-linear in $p$ and, by part~(i), agree when $c=0$; moreover both are
solutions of the first-order linear ODE $\partial_c F=-F'$ in $c$ with the same
initial value $F|_{c=0}=p$.  Indeed $p(x-c)$ obviously satisfies
$\partial_c[p(x-c)]=-p'(x-c)$; and for the convolution, since
$(x-c)^n=T^{(c)}x^n$ where $T^{(c)}=\exp(-c\,\partial_x)$ is the shift operator,
$p\boxp(x-c)^n=T^{(c)}p=p(x-c)$ by the same coefficient-matching argument as in
Proposition~\ref{prop:heatconv} (shift is the degree-one convolution operator).
Thus $p\boxp(x-c)^n=p(x-c)$.
Translation invariance of $\score,\Phi_n$ is then clear since they depend only
on root differences $\lambda_i-\lambda_j$, which are unchanged by $x\mapsto x-c$.

(iii) Replacing $a_k\mapsto s^k a_k$, $b_k\mapsto s^kb_k$ multiplies $c_k$ by
$s^k$ (every term in $c_k$ has $i+j=k$), which is the coefficient rescaling of
$(p\boxp q)_s$. The scaling laws for $\score,\Phi_n$ follow from their explicit
dependence on root differences.
\end{proof}

\section{The finite free Gaussian heat semigroup}\label{sec:heat}

Let $H_n$ be the monic probabilists' Hermite polynomial of degree $n$, the
unique monic degree-$n$ polynomial with $H_n''-xH_n'+nH_n=0$.  Its roots
$h_1,\dots,h_n$ satisfy, at each root, $H_n''(h_i)=h_iH_n'(h_i)$, hence
$2s_i(H_n)=H_n''(h_i)/H_n'(h_i)=h_i$, i.e.\
\[
\score(h)=\tfrac12 h. \tag{4.1}
\]
For $t>0$ set $G_t(x):=t^{n/2}H_n(x/\sqrt t)$ (roots $\sqrt t\,h_i$) and
$G_0:=x^n$.

\begin{proposition}[Heat flow $=$ convolution with $G_t$]\label{prop:heatconv}
For every monic $p$ of degree $n$ and $t\ge0$,
\[
p\boxp G_t \;=\; T_t\,p,\qquad T_t:=\exp\!\big(-\tfrac t2\,\partial_x^2\big)
=\sum_{m\ge0}\frac{(-t/2)^m}{m!}\,\partial_x^{2m},
\]
a finite sum on degree-$n$ polynomials.  Consequently $p_t:=p\boxp G_t$ solves
\[
\partial_t p_t=-\tfrac12\,\partial_x^2 p_t,\qquad p_0=p. \tag{H}
\]
\end{proposition}

\begin{proof}
Both sides are bilinear in the coefficients, so we match coefficients.  For
$p=\sum_k a_kx^{n-k}$,
\[
[x^{n-k}](T_tp)=\sum_{m\ge0}\frac{(-t/2)^m}{m!}\,\frac{(n-k+2m)!}{(n-k)!}\,a_{k-2m}.
\]
The Hermite coefficients are $b_{2m}=(-1)^m\,t^{m}\frac{n!}{(n-2m)!\,2^mm!}$,
$b_{\text{odd}}=0$ (these are exactly the coefficients of $G_t=T_t x^n$, which
follow from $H_n=T_1x^n$, the standard generating identity for Hermite
polynomials).  Substituting into
$c_k=\sum_{i+j=k}\frac{(n-i)!(n-j)!}{n!(n-k)!}a_ib_j$ with $j=2m$, $i=k-2m$
(so $n-i=n-k+2m$) gives
\[
\frac{(n-i)!(n-2m)!}{n!(n-k)!}\cdot(-1)^m t^m\frac{n!}{(n-2m)!\,2^mm!}
=\frac{(-t/2)^m}{m!}\,\frac{(n-k+2m)!}{(n-k)!},
\]
matching $[x^{n-k}](T_tp)$.  Hence $p\boxp G_t=T_tp$, and differentiating
$T_t=\exp(-\tfrac t2\partial^2)$ in $t$ gives (H).
\end{proof}

\begin{corollary}[Gaussian semigroup and its information]\label{cor:gauss}
$G_s\boxp G_t=G_{s+t}$ for $s,t\ge0$, and
$\Phi_n(G_t)=\dfrac{n(n-1)}{4t}$ for $t>0$.
\end{corollary}

\begin{proof}
$G_s\boxp G_t=T_tG_s=T_tT_sx^n=T_{s+t}x^n=G_{s+t}$ by
Proposition~\ref{prop:heatconv} and $T_tT_s=T_{s+t}$ (exponentials of a fixed
operator).  By (4.1), $\Phi_n(H_n)=\sum_i s_i(H_n)^2=\tfrac14\sum_i h_i^2$.
Writing $H_n=x^n+0\cdot x^{n-1}+c_2x^{n-2}+\cdots$ (the
$x^{n-1}$ coefficient vanishes because $\sum_ih_i=0$ by the recurrence
$H_n=xH_{n-1}-(n-1)H_{n-2}$, which also gives $c_2=-\binom n2$), Newton's
identity $\sum_i h_i^2=(\sum_i h_i)^2-2e_2(h)=-2c_2$ yields
$\sum_i h_i^2=n(n-1)$.  Therefore $\Phi_n(H_n)=\tfrac14 n(n-1)$, and by
Lemma~\ref{lem:cov}(iii), $\Phi_n(G_t)=t^{-1}\Phi_n(H_n)=\frac{n(n-1)}{4t}$.
\end{proof}

\section{Root dynamics and the exact de Bruijn identity}\label{sec:debruijn}

\begin{lemma}[Root velocity is the score]\label{lem:rootvel}
Let $p_t$ solve (H) with $p_0=p$ having simple roots, and let
$\lambda_i(t)$ be its (locally smooth, simple) roots.  Then
$\dot\lambda_i(t)=s_i(p_t)$.
\end{lemma}

\begin{proof}
Differentiating $p_t(\lambda_i(t))=0$:
$0=\partial_tp_t(\lambda_i)+\dot\lambda_i\,p_t'(\lambda_i)$, so by (H),
$\dot\lambda_i=\tfrac12 p_t''(\lambda_i)/p_t'(\lambda_i)$.  For a simple-rooted
$p$, $p'(x)=p(x)\sum_k(x-\lambda_k)^{-1}$ and
$p''(\lambda_i)/p'(\lambda_i)=2\sum_{j\neq i}(\lambda_i-\lambda_j)^{-1}=2s_i$
(standard).  Hence $\dot\lambda_i=s_i$.  Simplicity is preserved for small $t$
by the implicit function theorem since $p_t'(\lambda_i)\ne0$.
\end{proof}

\begin{theorem}[de Bruijn identity]\label{thm:debruijn}
Along the heat flow, for $p_t$ with simple roots,
\[
\frac{d}{dt}\,\Phi_n(p_t)=-\,I(p_t),\qquad
I(p_t)=\sum_{i\neq k}\frac{(s_i-s_k)^2}{(\lambda_i-\lambda_k)^2}=2\inner{s}{Ls}\ \ge\ 0 ,
\]
where $s=s(p_t)$.  In particular $\Phi_n(p_t)$ is nonincreasing in $t$.
\end{theorem}

\begin{proof}
With $\dot\lambda_i=s_i$ (Lemma~\ref{lem:rootvel}),
$\dot s_i=-\sum_{j\neq i}\frac{\dot\lambda_i-\dot\lambda_j}{(\lambda_i-\lambda_j)^2}
=-\sum_{j\neq i}\frac{s_i-s_j}{(\lambda_i-\lambda_j)^2}=-(Ls)_i$.
Hence, using $\Phi_n=\inner{s}{s}$,
\[
\frac{d}{dt}\Phi_n(p_t)=2\inner{s}{\dot s}=-2\inner{s}{Ls}
=-\sum_{i\neq j}\frac{(s_i-s_j)^2}{(\lambda_i-\lambda_j)^2}=-I(p_t)\le0,
\]
the last inequality by (2.1).
\end{proof}

\section{A sharp finite free Cram\'er--Rao / Stam lower bound}\label{sec:cr}

This is the key new ingredient.  Define, for $p$ with simple roots,
\[
g(p)\;:=\;\frac{I(p)}{\Phi_n(p)^2}\;=\;-\frac{d}{dt}\Big|_{t=0}\frac{1}{\Phi_n(p_t)},
\qquad p_t=p\boxp G_t,
\]
the chain-rule identity following from Theorem~\ref{thm:debruijn}
($\frac{d}{dt}\Phi_n^{-1}=-\Phi_n^{-2}\dot\Phi_n=I/\Phi_n^2$).  Note $g$ is
scale and translation invariant (Lemma~\ref{lem:cov}: $I\mapsto s^{-4}I$,
$\Phi_n^2\mapsto s^{-4}\Phi_n^2$).

\begin{theorem}[Sharp lower bound on $g$]\label{thm:gbound}
For every real-rooted monic $p$ of degree $n\ge2$ with simple roots,
\[
g(p)=\frac{I(p)}{\Phi_n(p)^2}\ \ge\ \frac{4}{n(n-1)},
\qquad\text{equivalently}\qquad
I(p)\ \ge\ \frac{4}{n(n-1)}\,\Phi_n(p)^2 ,
\]
with equality if and only if $p$ is a scaled, translated Hermite polynomial
$G_t$ (the finite free Gaussian).
\end{theorem}

\begin{proof}
Let $\lambda$ be the root vector and $s=\score(\lambda)$.  Recall from
Lemma~\ref{lem:Lid} that $L\lambda=s$, $\Phi_n=\inner{s}{s}$,
$I=2\inner{s}{Ls}$, and $\inner{\lambda}{L\lambda}=n(n-1)/2$.

Since $L$ is positive semidefinite, $\inner{u}{v}_L:=\inner{u}{Lv}$ is a
(possibly degenerate) inner product, and the Cauchy--Schwarz inequality holds:
\[
\inner{\lambda}{L s}^2\ \le\ \inner{\lambda}{L\lambda}\,\inner{s}{Ls}. \tag{6.1}
\]
The left side is $\inner{L\lambda}{s}^2=\inner{s}{s}^2=\Phi_n(p)^2$ (using
$L$ symmetric and $L\lambda=s$).  The right side is
$\frac{n(n-1)}{2}\cdot\inner{s}{Ls}=\frac{n(n-1)}{2}\cdot\frac{I(p)}{2}
=\frac{n(n-1)}{4}\,I(p)$.  Thus (6.1) reads
\[
\Phi_n(p)^2\ \le\ \frac{n(n-1)}{4}\,I(p),
\]
which is exactly $I(p)\ge\frac{4}{n(n-1)}\Phi_n(p)^2$, i.e.\
$g(p)\ge \frac{4}{n(n-1)}$.

Equality in Cauchy--Schwarz (6.1) holds iff $\lambda$ and $s$ are linearly
dependent modulo $\ker L=\operatorname{span}\{\mathbf 1\}$, i.e.\ iff there are
scalars $\alpha,\beta$ (not the trivial $\alpha=0$ unless $s\in\ker L$, which is
impossible for $n\ge2$ since $s\ne0$) with $s=\alpha\lambda+\beta\mathbf 1$.
By translation and scaling invariance (Lemma~\ref{lem:cov}) we may center and
scale so that $\beta=0$ and $\alpha=\tfrac12$, i.e.\ $\score(\lambda)=\tfrac12\lambda$;
by (4.1) and the fact that $H_n$ is the unique monic polynomial whose roots
satisfy $2\,\score=\lambda$ (this is the system $H_n''=xH_n'$ at the roots,
forcing $H_n''-xH_n'+nH_n=0$ since both sides are monic-degree-$n$ multiples of
$H_n$ vanishing at all roots), the equality case is exactly $p=G_t$ for some
$t>0$ (up to translation).
\end{proof}

\begin{remark}
Theorem~\ref{thm:gbound} is the finite free analogue of the classical
Cram\'er--Rao bound $I(X)\,\mathrm{Var}(X)\ge1$ rewritten in a scale-free form;
here the role of the ``variance'' is played by $n(n-1)/4$, the universal value
$\Phi_n(G_t)\cdot t$.  Its sharpness at the Hermite polynomial is the exact
analogue of equality for Gaussians.
\end{remark}

\section{The Gaussian case of $(\star)$: a complete proof for all $n$}\label{sec:gaussstar}

\begin{theorem}[Stam inequality against a Gaussian factor]\label{thm:gausscase}
Let $p$ be monic real-rooted of degree $n\ge2$ and let $G_t$ ($t>0$) be the
finite free Gaussian (any scaled, translated Hermite polynomial of degree $n$).
Then
\[
\frac{1}{\Phi_n(p\boxp G_t)}\ \ge\ \frac{1}{\Phi_n(p)}+\frac{1}{\Phi_n(G_t)},
\]
i.e.\ $(\star)$ holds with $q=G_t$.  Moreover the inequality is strict for every
$t>0$ unless $p$ is itself a (scaled, translated) Hermite polynomial, in which
case equality holds (consistently with $G_s\boxp G_t=G_{s+t}$).
\end{theorem}

\begin{proof}
By translation invariance (Lemma~\ref{lem:cov}(ii)) we may take $G_t$ centred.
For $u\ge0$ define
\[
\psi(u)\;:=\;\frac{1}{\Phi_n(p\boxp G_u)}-\frac{1}{\Phi_n(p)}-\frac{1}{\Phi_n(G_u)} .
\]
Here $p\boxp G_u$ is the heat flow $p_u$ of $p$ (Proposition~\ref{prop:heatconv}),
which has simple roots for all $u>0$ (heat flow strictly separates roots; in any
case $\Phi_n(p_u)<\infty$ for $u>0$, and at isolated coalescences both relevant
terms are continuous and the inequality is closed).  At $u=0$:
$p\boxp G_0=p$, and $1/\Phi_n(G_0)=0$ (since $G_0=x^n$ has $\Phi_n=\infty$),
so $\psi(0)=0$.

By Corollary~\ref{cor:gauss}, $\frac{1}{\Phi_n(G_u)}=\frac{4u}{n(n-1)}$, so
$\frac{d}{du}\frac{1}{\Phi_n(G_u)}=\frac{4}{n(n-1)}$.  By the chain rule and
Theorem~\ref{thm:debruijn},
$\frac{d}{du}\frac{1}{\Phi_n(p\boxp G_u)}=g(p\boxp G_u)$.  Hence
\[
\psi'(u)\;=\;g(p\boxp G_u)\;-\;\frac{4}{n(n-1)}\;\ge\;0
\]
for all $u>0$, by Theorem~\ref{thm:gbound} applied to $p\boxp G_u$.  Since
$\psi(0)=0$ and $\psi$ is nondecreasing, $\psi(t)\ge0$, which is the claimed
inequality.

For the equality discussion: $\psi(t)=0$ forces $\psi'\equiv0$ on $[0,t]$,
hence $g(p\boxp G_u)=\frac{4}{n(n-1)}$ for all $u\in(0,t)$; by the equality
case of Theorem~\ref{thm:gbound} this means $p\boxp G_u=G_{c(u)}$ is Hermite
for all such $u$, which (taking $u\to0^+$, using continuity of roots) forces
$p=G_{c(0)}$ to be Hermite.  Conversely if $p=G_s$ then
$p\boxp G_t=G_{s+t}$ and equality holds by Corollary~\ref{cor:gauss}:
$\frac1{\Phi_n(G_{s+t})}=\frac{4(s+t)}{n(n-1)}=\frac{4s}{n(n-1)}+\frac{4t}{n(n-1)}
=\frac1{\Phi_n(G_s)}+\frac1{\Phi_n(G_t)}$.
\end{proof}

\section{The case $n=2$}\label{sec:n2}

\begin{proposition}\label{prop:n2}
For $n=2$ and any real-rooted monic $p,q$, equality holds in $(\star)$.
\end{proposition}

\begin{proof}
By translation invariance take roots of $p$ to be $\{0,d_p\}$ and of $q$ to be
$\{0,d_q\}$.  Then $a=(1,-d_p,0)$, $b=(1,-d_q,0)$, and
\[
c_0=1,\quad c_1=-d_p-d_q,\quad
c_2=\tfrac12(a_0b_2+a_2b_0)+a_1b_1=\tfrac12 d_pd_q .
\]
The roots of $x^2-(d_p+d_q)x+\tfrac12 d_pd_q$ have squared gap
$\Delta^2=(d_p+d_q)^2-2d_pd_q=d_p^2+d_q^2$.  For $n=2$,
$\Phi_2=2/\Delta^2$ (Lemma~\ref{lem:phiform}), so
$\Phi_2(p)=2/d_p^2$, $\Phi_2(q)=2/d_q^2$,
$\Phi_2(p\boxplus_2q)=2/(d_p^2+d_q^2)$, whence
$1/\Phi_2(p\boxplus_2q)=(d_p^2+d_q^2)/2=1/\Phi_2(p)+1/\Phi_2(q)$.
\end{proof}

\section{Degenerate cases}\label{sec:degenerate}

If $p$ has a repeated root then $\Phi_n(p)=\infty$, $1/\Phi_n(p)=0$, and
$(\star)$ reduces to $1/\Phi_n(p\boxp q)\ge1/\Phi_n(q)$.  The coefficients of
$p\boxp q$ are polynomial (continuous) functions of the roots of $p,q$, the
roots of $p\boxp q$ depend continuously on those coefficients, and $1/\Phi_n$
is continuous on the simple-root locus and extends by $0$ to the repeated-root
locus (Lemma~\ref{lem:phiform}).  Thus it suffices to prove the relevant
inequality for simple roots and pass to the limit; this applies in particular
to Theorems~\ref{thm:gausscase} (taking limits in $p$) and to $(\star)$ in
general.

\section{The general case: an honest status and the exact obstruction}
\label{sec:gap}

We now give a careful account of what remains, correcting a tempting but
\emph{incorrect} reduction.

\subsection*{Equivalent reformulation}
By a one-variable optimization, $(\star)$ for fixed $p,q$ with $r=p\boxp q$ is
\emph{equivalent} to
\[
\Phi_n(r)\ \le\ \min_{\lambda\in[0,1]}
\big(\lambda^2\Phi_n(p)+(1-\lambda)^2\Phi_n(q)\big)
=\frac{\Phi_n(p)\Phi_n(q)}{\Phi_n(p)+\Phi_n(q)} ,
\tag{$\star'$}
\]
the minimum being attained at $\lambda^\star=\Phi_n(q)/(\Phi_n(p)+\Phi_n(q))$.
Indeed $\frac{\Phi_n(p)\Phi_n(q)}{\Phi_n(p)+\Phi_n(q)}$ is the reciprocal of
$\frac1{\Phi_n(p)}+\frac1{\Phi_n(q)}$.  In the classical Stam proof one
realizes $\score(r)$ as the value of a conditional expectation,
$\score(r)=\mathbb E[\lambda^\star\score(p)+(1-\lambda^\star)\score(q)\mid\cdot]$,
and ($\star'$) follows from the $L^2$-contractivity of conditional expectation.
The finite free analogue of this statement -- a \emph{score-projection lemma}
expressing $\score(r)$ as a norm-nonincreasing combination of $\score(p)$ and
$\score(q)$ -- is the exact missing ingredient.

\subsection*{Why the coupled-heat-flow ``infinitesimal Stam'' route does NOT
close the general case}
A natural attempt mimics Blachman: run the heat flow on both factors with a
schedule $(a(t),b(t))$, set $P=p\boxp G_a$, $Q=q\boxp G_b$, $R=P\boxp Q$, and
use $R\boxp G_{a+b}'$-type identities.  Using
$r\boxp G_{a+b}=(p\boxp G_a)\boxp(q\boxp G_b)$ (Corollary~\ref{cor:gauss}
with commutativity/associativity) and the de~Bruijn identity, the gap
$D(t)=\frac1{\Phi_n(R)}-\frac1{\Phi_n(P)}-\frac1{\Phi_n(Q)}$ has derivative
\[
D'(t)=(\dot a+\dot b)\,g(R)-\dot a\,g(P)-\dot b\,g(Q). \tag{D}
\]
To make $D$ monotone one would need, at every reachable state, nonnegative
speeds $(\dot a,\dot b)$, not both zero, with
\[
(\dot a+\dot b)\,g(R)\ \le\ \dot a\,g(P)+\dot b\,g(Q). \tag{S}
\]
\emph{This is impossible in general.}  The linear inequality (S) in
$(\dot a,\dot b)\ge0$ is feasible (with a nonzero solution) if and only if
$\min\{g(P),g(Q)\}\le g(R)$ fails, i.e.\ it is \emph{infeasible} precisely when
\[
g(R)\ >\ \max\{g(P),g(Q)\}. \tag{$\dagger$}
\]
We verified numerically that ($\dagger$) occurs at genuinely reachable states,
including at the initial time $t=0$ with $P=p$, $Q=q$: e.g.\ for $n=4$ there are
real-rooted $p,q$ with
$g(p)\approx0.615,\ g(q)\approx0.572,\ g(p\boxp q)\approx0.829$,
so $g(p\boxp q)>\max\{g(p),g(q)\}$ and (S) has \emph{no} admissible solution at
$t=0$.  (Many such instances were found across $3\le n\le5$.)  Hence no choice
of schedule can make $D$ monotone from the start, and the de~Bruijn coupling
\emph{alone} cannot prove $(\star)$.  This corrects the previously circulated
``reduction to (S)'', whose hypothesis is false at reachable states and which
is therefore vacuous as a proof strategy.

The phenomenon is consistent with the classical theory: there, too, the naive
pointwise inequality $2g(R)\le g(P)+g(Q)$ is false; the classical proof does
\emph{not} use a pointwise speed inequality but the score-projection
(conditional-expectation) identity directly.  The finite free obstruction
($\dagger$) is exactly the statement that no such projection is visible at the
level of the scalar quantities $g$ alone; one genuinely needs the vector-level
score-projection lemma.

\subsection*{Summary of the status of $(\star)$}
\begin{itemize}
\item $(\star)$ is \emph{proved unconditionally} for $n=2$ (equality,
Prop.~\ref{prop:n2}) and, for all $n$, \emph{whenever one factor is a finite
free Gaussian} $G_t$ (Theorem~\ref{thm:gausscase}), via the sharp bound
$g\ge\frac4{n(n-1)}$ (Theorem~\ref{thm:gbound}).
\item Extensive numerical search ($>2\times10^4$ random instances, $2\le n\le5$)
found \emph{no} violation of $(\star)$, of ($\star'$), or of the analogous
multiplicative form for the dissipation $I$; equality is approached exactly as
one factor tends to a Gaussian, matching Theorem~\ref{thm:gausscase}.
\item The general case is reduced to a single \emph{finite free
score-projection lemma}: $\score(p\boxp q)$ is a norm-nonincreasing affine
combination of $\score(p)$ and $\score(q)$ (the optimal weight being
$\lambda^\star$ above).  This remains open.
\end{itemize}

\section{Summary}

We proved rigorously and unconditionally: the closed form for $\Phi_n$ and the
score-Laplacian identities (Lemmas~\ref{lem:phiform},~\ref{lem:Lid});
covariance and the identity element (Lemma~\ref{lem:cov}); that finite free
convolution with the Hermite Gaussian $G_t$ is the heat flow
$\partial_tp=-\tfrac12 p''$ and the Gaussian semigroup with
$\Phi_n(G_t)=\frac{n(n-1)}{4t}$ (Prop.~\ref{prop:heatconv},
Cor.~\ref{cor:gauss}); the identification of root velocity with the score and
the exact de~Bruijn identity with nonnegative dissipation
(Lemma~\ref{lem:rootvel}, Theorem~\ref{thm:debruijn}); the \emph{sharp finite
free Cram\'er--Rao/Stam bound} $I\ge\frac{4}{n(n-1)}\Phi_n^2$ with equality
exactly at Hermite (Theorem~\ref{thm:gbound}); and -- as the main consequence
-- the \emph{complete} proof of $(\star)$ against a Gaussian factor for all $n$
(Theorem~\ref{thm:gausscase}), together with $n=2$ (Prop.~\ref{prop:n2}).  We
also corrected the false ``infinitesimal Stam'' reduction by exhibiting
reachable states where its hypothesis is infeasible, and isolated the exact
remaining open ingredient, a finite free score-projection lemma.
\end{document}
